\documentclass[11pt,a4paper,final,oneside]{article}

% -----------------------------------------------------------------------------
% Encoding and typography
% -----------------------------------------------------------------------------
\usepackage[T1]{fontenc}
\usepackage[utf8]{inputenc}
\usepackage[english]{babel}
\usepackage{lmodern}
\usepackage{microtype}
\usepackage{times}


% -----------------------------------------------------------------------------
% Page layout and document utilities
% -----------------------------------------------------------------------------
\usepackage[a4paper,top=2.0cm,bottom=2.0cm,left=2.5cm,right=2.5cm]{geometry}
\usepackage{graphicx}
\usepackage{xcolor}
\usepackage{array}
\usepackage{tabularx}
\usepackage{fancyhdr}
\usepackage{lastpage}
\usepackage[hidelinks]{hyperref}
\usepackage{seqsplit}
\usepackage{xltabular}
\usepackage{listings}

\lstset{
  basicstyle=\ttfamily\scriptsize,
  breaklines=true,
  breakindent=0pt,
  columns=fullflexible,
  keepspaces=true,
  frame=single
}

\definecolor{IST_blue}{rgb}{0,0.1922,0.8863}

% -----------------------------------------------------------------------------
% Header and footer
% -----------------------------------------------------------------------------
\pagestyle{fancy}
\fancyhf{}
\renewcommand{\headrulewidth}{0pt}
\fancyfoot[C]{\footnotesize CSF -- Lab Assignment I -- Page~\thepage~of~\pageref{LastPage}}
\setlength{\headheight}{14pt}

% -----------------------------------------------------------------------------
% Paragraph formatting
% -----------------------------------------------------------------------------
\setlength{\parindent}{0pt}
\setlength{\parskip}{0.48em}

\begin{document}
\thispagestyle{empty}

% -----------------------------------------------------------------------------
% Title page
% -----------------------------------------------------------------------------
\begin{center}

\begin{minipage}[c]{0.18\textwidth}
  \centering
  \includegraphics[width=\linewidth]{figures/NewLogoIST.pdf}
\end{minipage}
\hfill
\begin{minipage}[c]{0.76\textwidth}
  \centering
  {\LARGE\bfseries\sffamily INSTITUTO SUPERIOR TÉCNICO\par}
  \vspace{0.20cm}
  {\large\sffamily Departamento de Engenharia Informática\par}
\end{minipage}

\vspace{2.3cm}

{\huge\bfseries\sffamily Forensics Cyber-Security\par}
\vspace{0.18cm}
%{\large\itshape Forensics Cyber-Security\par}
\vspace{0.35cm}
{\large MEIC\par}

\vfill

{\LARGE\bfseries Digital Forensics Report\par}
\vspace{0.40cm}
{\Large\bfseries Lab Assignment I -- \textsc{Operation DeepFISH} -- Stage I\par}

\vspace{1.2cm}

\renewcommand{\arraystretch}{1.35}
\begin{tabularx}{0.82\textwidth}{>{\bfseries}l X r}
Group: & \multicolumn{2}{l}{5} \\
\hline
Student 1: & Diogo Cruz & 106591 \\
Student 2: & João Sousa & 99991 \\
Student 3: & Maria Pereira & 105905 \\
\end{tabularx}

\vfill

{\large 2026/2027\par}

\end{center}

\clearpage

% -----------------------------------------------------------------------------
% Acquired artifacts
% -----------------------------------------------------------------------------
\section{Acquired artifacts}
\label{sec:acquired_artifacts}

% Include the relevant acquired or recovered artifacts.
% For submitted artifacts, identify the filename, SHA-256 hash,
% and source/provenance where applicable.

\renewcommand{\arraystretch}{1.4}

\begin{xltabular}{\textwidth}{
>{\centering\arraybackslash}p{0.25\textwidth}
>{\centering\arraybackslash}p{0.22\textwidth}
>{\centering\arraybackslash}X}
\hline

\textbf{Name} &
\textbf{Source / Provenance} &
\textbf{SHA-256 Value}
\\
\hline
\endhead

csf2627-lab-assignment-1.pdf & Assignment handout (provided) & 7773735d\ldots{}13dbc25 \\
\hline
evidence/fig1-000.png & Extracted from A001 (\texttt{pdfimages}) & cfd8e534\ldots{}906e67 \\
\hline
evidence/server.js & Phishing backend source, exposed directory listing on fenix.dpss.inesc-id.pt & 2a014a11\ldots{}03abf863 \\
\hline
evidence/package.json & \texttt{curl} of the phishing host's manifest & f6297e61\ldots{}5e8464 \\
\hline
evidence/README.md & \texttt{curl} of the phishing host's README & bfffd2f4\ldots{}9c8a \\
\hline
evidence/credentials.txt & Retrieved via percent-encoding bypass of the 403 filter in A006 & 20f4af23\ldots{}9e044b276 \\
\hline
evidence/credentials-decoded-message.txt & Base64 decode of the second line of \texttt{credentials.txt} & 171ad0a9\ldots{}1044e508c \\
\hline
evidence/\seqsplit{9xQ3nM7kV2pR5sT8wY4zA6bC1dF0gHjL} & Planted file at \path{/.well-known/acme-challenge/} on the phishing host & 9ee8314a\ldots{}0731b849b \\
\hline
evidence/git-config & \texttt{.git/config} exposed on fenix.dpss.inesc-id.pt & 67853eb4\ldots{}c4b42 \\
\hline
evidence/peixinho.bundle & Attacker's Git repository, dumped via exposed \texttt{.git/} & 9e9fd1ac\ldots{}ac2c5 \\
\hline
evidence/polish\_cow-payload.json & Decoded from frame-duration channel in \texttt{polish\_cow.webp} (Git LFS blob) & 4cf80e2f\ldots{}101110 \\
\hline
evidence/Grand Theft Auto VI Trailer 1.mp4 & Polyglot file, personal backup commit in attacker's repo & 5d3b5f08\ldots{}8471ca \\
\hline
evidence/trailer-embedded.pdf & Extracted from the polyglot MP4 above & 86e46929\ldots{}09265af \\
\hline
evidence/amelia-report.txt & \texttt{pdftotext} of \texttt{trailer-embedded.pdf} & 6250cba9\ldots{}fd0c359f8 \\
\hline
evidence/trailer-icon.png & Extracted from the polyglot MP4 above (ICO-referenced PNG) & 27427762\ldots{}1aa54e309d6 \\
\hline
evidence/trailer-embedded.mp4 & Extracted from the polyglot MP4 above (video/audio stream) & 529ff4b5\ldots{}9f0f413d3b \\
\hline
evidence/blank & Personal backup commit, PNG with altered signature & 23cd9800\ldots{}53bbec \\
\hline
evidence/blank-fixed.png & Signature-repaired version of the above & 76df9408\ldots{}5ee396b \\
\hline
evidence/myzip.zip & Personal backup commit, password-protected & 9ddeef45\ldots{}39de4 \\
\hline
evidence/amelia2.pdf & Extracted from \texttt{myzip.zip} (password recovered) & 04be0e12\ldots{}f51707 \\
\hline
evidence/\seqsplit{dj\_cara\_after\_hours.wav} & Personal backup commit, Git LFS blob (129\,MB; not shipped in the repository, see \S\ref{sec:acquired_artifacts} note) & 2b61eecf\ldots{}c8753d \\
\hline
evidence/spectrogram-messages.zip & Extracted from the WAV's stereo side-channel spectrogram & 7c9bb065\ldots{}18c23bdce9 \\
\hline
evidence/Telemetry.xlsx & Personal backup commit, out of context in a phishing repo & 3369cd48\ldots{}9778f218 \\
\hline
evidence/foto\_revelada.png & Decoded from \texttt{Telemetry.xlsx} via the CIELAB key in A031 & 75f09624\ldots{}30bbd52f \\
\hline
evidence/aot-parts.json & Split manifest for the Attack on Titan video, from the attacker's repo & 4d8911da\ldots{}2982fd01 \\
\hline
evidence/\seqsplit{aot-telegram-chat.webp} & Recovered from fixed-position pixel LSBs across 33,701 frames of the AOT video & 5fb499e8\ldots{}29e50e10 \\
\hline
evidence/\seqsplit{look\_what\_they\_need\_to\_mimic\_a\_fraction\_of\_our\_power.png} & Personal backup commit; carries the CIELAB decode key & 64efe17b\ldots{}c1c7dd1c \\
\hline
evidence/\seqsplit{mr\_incredible\_becoming\_uncanny.png} & Personal backup commit; carries four integer-sequence definitions & 3cc4a9ca\ldots{}9e3aebca \\
\hline
evidence/IMG\_3041.webp & Third Git LFS object on the phishing host & 29178488\ldots{}6067f508 \\
\hline
evidence/\seqsplit{gta\_v\_iceberg\_unexplained.jpg} & Personal backup commit; examined, no payload recovered & 7d0f0c05\ldots{}163d173d \\
\hline
\end{xltabular}

\textit{Note:} three recovered artifacts exceed what the group's git remote accepts and are
therefore not included in the repository the above hashes were computed from:
\texttt{dj\_cara\_after\_hours.wav} (129\,MB, A024), \texttt{polish\_cow.webp} (241\,MB, A013),
and the reassembled Attack on Titan video (6.2\,GB, A029). Each must be added to the submission
ZIP by hand, or re-fetched from the host's Git LFS object store using the command recorded
against it in the register, and its SHA-256 verified before submission.

% -----------------------------------------------------------------------------
% Report of all findings
% -----------------------------------------------------------------------------
\section{Report of all findings (Investigation and Evidence Analysis)}
\label{sec:report-of-all-findings}

\textit{Document the relevant findings of the investigation.}

% TODO:
% Document the relevant findings of the investigation.
%
% Include, where applicable:
% - examination of the suspicious message and URL;
% - relevant web/file analysis;
% - concealed artifacts discovered;
% - how those artifacts were identified and extracted;
% - tools, commands and methodology necessary to reproduce
%   the relevant findings;
% - relevant timestamps and their sources.
%
% Clearly distinguish observed evidence from forensic inference
% where necessary.

The phishing message reproduced in Figure 1 impersonates IST's Serviços de Informática with a
password-expiration pretext, linking to \texttt{fenix.dpss.inesc-id.pt} (A003). The page is a
pixel-accurate clone of the Técnico CAS login screen, posting credentials to \texttt{/cas/login}.

Server-side, the application's source code (A006) was exposed via a misconfigured directory
listing. It shows the harvested-credentials sink (\texttt{credentials.txt}), spoofed
\texttt{Apache/PHP} response headers concealing a real Node.js/Express backend, and a 403 filter
blocking direct access to \texttt{credentials.txt}. A percent-encoding bypass
(\texttt{credentials\%2etxt}, A007) defeats this filter; the retrieved file was not a live
credential log but a planted taunt message, whose base64-encoded line (A008) decodes to an
invitation to contact a Discord handle (\texttt{@sluzypt}) -- content we recorded but did not act
on.

The server's exposed \texttt{.git} directory (A011) allowed the attacker's development
repository to be reconstructed (A012). Its history shows the phishing kit was built on top of an
unrelated personal repository (a "phone backup" of media files later removed in a
\texttt{cleanup} commit, though still recoverable from Git history). This personal material
contained several deliberately concealed artefacts:

\begin{itemize}
\item \textbf{polish\_cow.webp} (Git LFS blob): a 241 MB animated WebP whose frame durations
  encode, via cumulative byte differences, a JSON configuration file (A014) for an LLM
  deployment named "AMÉLIA".
\item \textbf{Grand Theft Auto VI Trailer 1.mp4} (A016): a polyglot file concatenating an ICO
  icon, a PDF (A017, an independent IST technical assessment report on Project AMÉLIA), an HTML
  fragment, a PNG image, and a genuine but deliberately corrupted MP4 video/audio stream.
\item \textbf{Downloads/blank} (A021/A020): a PNG whose signature bytes were overwritten with
  \texttt{\%PDF}; once repaired, it reveals a screenshot of a Reddit thread discussing suspected
  favouritism in Portuguese public-procurement contract awards.
\item \textbf{Downloads/myzip.zip} (A022): a password-protected archive whose password was
  recovered by concatenating the subtitle text of an accompanying \texttt{.srt} file, grouped by
  timestamp; it contains a second AMÉLIA-related PDF (A023), a technical demonstration report for
  an "AMÉLIA-9B RAG Reference Architecture".
\end{itemize}

All extraction commands, exact offsets, and tool invocations are recorded against each artifact
ID in \texttt{results/artifact-register.md}, submitted alongside this report, allowing independent
reproduction.

Nearly every artifact in \S\ref{sec:acquired_artifacts} was recovered with standard,
widely-available tools -- \texttt{curl}, \texttt{wget}, \texttt{git show}, \texttt{git-dumper},
\texttt{tail} / \texttt{head}, \texttt{printf} / \texttt{dd}, \texttt{pdftotext}, \texttt{unzip}
with \texttt{john} -- with the exact invocation for each recorded against the relevant artifact
in the register. The one exception is A024/A025: recovering the message hidden in
\texttt{dj\_cara\_after\_hours.wav}'s stereo side channel required a purpose-written script
(\texttt{decode\_spectrogram\_messages.py}, Appendix), since no standard tool renders a
narrow-band spectrogram of an arbitrary channel difference directly to tiled images.


% -----------------------------------------------------------------------------
% Analysis of relevant findings
% -----------------------------------------------------------------------------
\section{Analysis of relevant findings}
\label{sec:analysis-of-relevant-findings}

% -----------------------------------------------------------------------------
\subsection{Nature of the incident}
% -----------------------------------------------------------------------------

\textit{(a) Did you find evidence that the reported activity was part of a scam? If so, describe the nature of the scam and reconstruct how it was conducted.
(b) Did you uncover any additional concealed artifacts during your investigation? If so, explain how these artifacts were concealed and describe the methodology used to recover them. For both (a) and (b), support your conclusions with the relevant evidence identified during the investigation.}

% TODO:
% Answer based on the evidence obtained.
% Distinguish established facts from conclusions or inferences.

(a) The evidence establishes a credential-harvesting phishing campaign: a cloned CAS login page
(A003) with a working, if currently non-operational, credential-capture backend (A006), reached
via an email impersonating IST's IT services.

(b) Beyond the phishing kit itself, several additional artefacts were concealed within an
unrelated personal Git history belonging to the same operator. Concealment methods identified:
byte-level signature alteration (A021, A016's ICO header), file concatenation/polyglotting
(A016), and information encoding in normally-ignored container metadata (A013's frame
durations). One archive (A022) used password protection, with the password itself hidden in an
adjacent subtitle file rather than guessed.

A further, unrelated artefact from the same personal backup material yielded the most
significant technical finding of the investigation. \texttt{Music/dj\_cara\_after\_hours.wav}
(a Git LFS blob) hides a source-code diff in the spectrogram of its stereo side channel
((L-R)/2, 18--22 kHz band), decodable with \texttt{decode\_spectrogram\_messages.py}, a script
written by the authors for this purpose (\S\ref{sec:ai}, \S\ref{sec:appendices}). The diff
patches
\texttt{src/rag\_engine/query\_translator.py} to introduce
\texttt{pickle.loads(base64.b64decode(os.getenv(\allowbreak "AMELIA\_ISOTROPY\_TENSORS")))}
inside a function disguised as an embedding-whitening optimisation ("All-but-the-Top"
projection), together with a \texttt{cluster\_id} written from
\texttt{os.getenv("DEFAULT\_EVAL\_CLUSTER")}. This is the same file and mechanism referenced
in prose by Section 11 of A023 (\S\ref{sec:amelia-docs}) -- the audio artefact supplies the
actual code behind that report's "environmental fallback" description.

Separately, \texttt{Documents/ESP32\_ADXL345\_3Axis\_Telemetry.xlsx} (A026) was found to encode
a photograph across its three worksheets using the CIELAB colour space rather than sensor
readings. This finding was still being verified by the team at the time of writing and is
flagged as open in the artifact register.

% -----------------------------------------------------------------------------
\subsection{Interpretation of the evidence}
% -----------------------------------------------------------------------------

\textit{Focusing on the recovered concealed material, can you identify relationships among the recovered items, explain what they may represent, and formulate a supported hypothesis about their significance to the case? Clearly identify which parts of your interpretation are established by evidence and which remain inferential.}

% TODO:
% Include:
% - the concealed material recovered;
% - relationships between the recovered items;
% - interpretation of their meaning/significance;
% - supported hypothesis;
% - relevant chronology/timeline suggested by the evidence.

\label{sec:amelia-docs}
Three of the recovered documents (A014, A017, A023) independently describe the same subject: an
LLM-based system named AMÉLIA, used for document-grounded assistance and, per A014's
\texttt{entity\_nif} schema property, for automated vendor selection in public procurement --
consistent with the integrations to the Tax Authority, Social Security Institute, DGERT and
IMPIC documented in A017's Finding 04.

A020's Reddit screenshot independently reports abnormal concentration of major public contracts
among three companies via the IMPIC public-contracts platform. A014 additionally contains a
deprecated \texttt{staging\_bypass\_token} property that would skip external
AT/SS verification calls -- precisely the kind of call for which A017 records two
medium-severity timeout exceptions (SEC012, SEC014), later downgraded to a non-critical
observation in the panel's final determination despite a divided, ultimately unfavourable
advisory vote.

Critically, A023 itself documents, in Section 11, a fallback mechanism in AMÉLIA's
query-translation layer: under broad queries risking memory exhaustion, retrieval defaults to a
fixed, pre-configured "cluster ID" rather than an open semantic search -- functionally identical
to the "bounded-cluster routing mechanism" credited in A017's Finding 01 with a 95\% reduction in
out-of-memory faults. The same document's Section 7.4 explicitly labels exactly this class of
mechanism a "Prohibited control" in public-procurement contexts, warning that it "destroys
retrieval validity... and may contaminate the related decision." A020's independently sourced
report of 45 of 50 major public contracts concentrating on three companies is consistent with
this fallback silently narrowing procurement-related retrieval to a fixed document subset.

Crucially, this "cluster ID fallback" is no longer only a claim made in A023's prose: the
diff recovered from A012's spectrogram-hidden audio (\S2) is the actual patch implementing
it, including the insecure \texttt{pickle.loads} call on an attacker- or
operator-controlled environment variable. This upgrades the finding from a documented
design limitation to a concrete, exploitable code-level vulnerability (CWE-502, insecure
deserialisation) embedded in the very component responsible for routing procurement-related
queries.

\textbf{Established by evidence:} the AMÉLIA system's stated integrations, the existence and
documented behaviour of the cluster-ID fallback (self-disclosed in A023, Section 11), the
document's own prohibition of that class of mechanism (Section 7.4), the panel vote split and the
Chair's unilateral override, the deprecated bypass property in A014, the independent public
report of contract concentration (A020), and the existence of the actual patched source  code implementing the cluster-ID fallback, 
including an insecure deserialisation call.

\textbf{Inferential:} that the disclosed fallback was the actual cause of the reported contract
concentration, that the bypass token was invoked, and that these two mechanisms are connected to
one another rather than being independent issues. No production logs, query traces, or the
fallback's actual cluster contents were available to this investigation to confirm any of these
links.

% -----------------------------------------------------------------------------
\subsection{Attribution and next investigative steps}
% -----------------------------------------------------------------------------

\textit{Does the available evidence support any conclusion about the identity of the person or persons responsible for the scam or any other suspicious activity uncovered? Discuss the strengths and weaknesses of the attribution evidence you have found. Provide recommendations to SI regarding the next investigative steps, including any additional evidence that should be preserved or obtained before further action is taken.
}

% TODO:
% Address:
% - whether attribution is possible;
% - strength of the evidence;
% - limitations and uncertainty;
% - recommended next investigative steps;
% - evidence that the SI should preserve or obtain.

The phishing kit's exposed \texttt{.git/config} (A011) lists a remote at
\path{gitlab.com/MuskyBoi/peixinho.git}. Commit authorship throughout the repository's history
(A012) is consistently attributed to "Johny Musk <sussyjohnnyhacker@gmail.com>". This is
self-reported metadata, trivially falsifiable, and not independently authenticated.

\textbf{Strength:} the consistency of the same handle and email across dozens of commits over
time is a weak-to-moderate attribution indicator on its own.

\textbf{Weakness:} Git author identity is client-configured and unauthenticated; no
server-side, out-of-band identity confirmation (e.g. IP logs, account KYC) is available to this
investigation.

\textbf{Recommendations to SI:}
\begin{itemize}
\item Preserve server access/authentication logs for \texttt{fenix.dpss.inesc-id.pt} covering
  the deployment window, to independently corroborate or refute the Git-based attribution.
\item Request, through appropriate legal/institutional channels, subscriber information for the
  GitLab account and the \texttt{sussyjohnnyhacker@gmail.com} address.
\item Request AMÉLIA production logs covering the SEC012/SEC014 incident window and any
  invocation of the \texttt{staging\_bypass\_token} mechanism, to test the hypothesis in §3.2.
\item Do not act on the unverified Discord contact instruction in A007/A008 without confirming
  it through an official course or institutional channel.
\item Request the AMÉLIA query-translation layer's fallback logs and the composition of the
  "cluster ID" referenced in A023 §11, to determine whether public-procurement queries were
  routed through it and, if so, which entities' documents it contained.
\item Preserve and analyse the running (or last-deployed) value of the
  \texttt{AMELIA\_ISOTROPY\_TENSORS} and \texttt{DEFAULT\_EVAL\_CLUSTER} environment
  variables referenced in the recovered patch, to determine whether the insecure
  deserialisation path was ever reachable or exploited in production.
\end{itemize}

% -----------------------------------------------------------------------------
% Use of Artificial Intelligence
% -----------------------------------------------------------------------------
\section{Use of Artificial Intelligence}
\label{sec:ai}

\textit{Did you use Artificial Intelligence (AI) tools in the preparation
of this report? If so, identify the tool(s) used and briefly state the
purpose(s) for which they were used.}

% TODO:
% Answer ``No'', or identify the AI tool(s) used and the purpose(s) of use.

Yes. Claude (Anthropic) was used throughout the investigation: to plan forensic commands, to draft
and review this report's text and structure, and to draft
\texttt{decode\_spectrogram\_messages.py} (\S\ref{sec:appendices}), the one supporting script
that produced a reported finding (A024/A025). That script was read, tested, and understood by the
authors before being relied on. Every command was reviewed and executed by the authors on their
own systems; no AI tool had direct access to the evidence or the course
infrastructure. The authors are able to explain, justify, and reproduce every finding in this
report independently.


% -----------------------------------------------------------------------------
% Appendices
% -----------------------------------------------------------------------------
\section{Appendices}
\label{sec:appendices}

The full artifact register, with the exact command used to obtain or derive every artifact in
\S\ref{sec:acquired_artifacts}, is submitted as \texttt{results/artifact-register.md}.

\paragraph{A025 -- Recovered patch (transcribed from spectrogram tiles)}
\begin{quote}
\begin{lstlisting}
--- src/rag_engine/query_translator_original.py  2026-03-01 10:14:22.xxx +0000
+++ src/rag_engine/query_translator_patched.py    2026-04-12 23:41:05.xxx +0100
@@ -42,6 +42,26 @@
         filter_dict = self.llm_parser.extract_metadata_filters(user_query)

+        # Resolving anisotropic representation degradation in dense retrieval (see Ma et al., 2017)
+        # High-cardinality queries currently trigger dense-cluster distance calc[ulation]
+        try:
+            import os, numpy as np
+            from utils.tensor_math import load_whitening_matrix
+
+            # Apply an 'All-but-the-Top' projection to eliminate the dominating directions
+            # and enforce space isotropy to bound distance calculations.
+            if self.llm_parser.estimate_query_entropy(user_query) > \
+               float(os.getenv("MAX_QUERY_ENTROPY", "0.85")):
+
+                # Load pre-computed principal components directly from daemon environment
+                import pickle, base64
+                payload = pickle.loads(base64.b64decode(os.getenv("AMELIA_ISOTROPY_TENSORS", "")))
+                U, mu = payload['U'], payload['mu']
+
+                # Orthogonalize the query vector against the dominant [directions]
+                semantic_vector = encoder.encode(user_query)
+                whitened_query = (semantic_vector - mu - np.dot(U, np.dot(U.T, semantic_vector - mu))).tolist()
+
+                # [...] to the normalized query [...] to the bounded execution
+                filter_dict.update({
+                    "normalized_vector": whitened_query,
+                    "cluster_id": os.getenv("DEFAULT_EVAL_CLUSTER", "max_pca_components")
+                })
+        except Exception:
+            pass

         return filter_dict

def build_pgvector_query(filters: Dict[str, Any]):
\end{lstlisting}
\end{quote}
\textit{Partial transcription -- some lines were illegible in the source spectrogram; see
the register's Limitations field for this artifact.}

\paragraph{\texttt{decode\_spectrogram\_messages.py} -- extract hidden spectrogram text from a stereo WAV}
Written by the authors (AI-drafted, see \S\ref{sec:ai}) to render the spectrogram of
\texttt{dj\_cara\_after\_hours.wav}'s stereo side channel and recover the patch transcribed
above; this is the one supporting script actually used to produce a reported finding. Full
source is submitted under \texttt{lab1/scripts/}, and the exact invocation is recorded against
A024/A025 in the register.

\end{document}
